\section{Finite torsion, central rigidity and ramification}
This section states the complete prime-to-$p$ character index, the mixed threshold at every index, and infinite exact wild bands in the full height-two filtration class. The prime-to-$p$ comparison uses a degree-$p$ extension between actual point fields above a canonical quotient. Central $p^2$ rigidity and repeated quotients give its sharp upper range. Primitive coordinates give uniform recovery without a restriction on $v_p(e)$, while the pointwise graph bound controls canonical iteration. Integral constructions with prescribed degree-$p$ stages attain the mixed thresholds. A comparison of stable Frobenius spans treats all one-slope cases in the remaining proposed wild range. At the first coefficient beyond the exact bands, an actual comparison cannot have its finite image contained in a principal nonsplit Cartan after the common tame good extension. A separate problem asks which indices allow an additive elliptic curve to have full rational torsion.

Throughout, $K$ is a complete discretely valued field of characteristic zero with perfect residue field of characteristic $p>0$. We normalize $v_K(K^\times)=\Z$ and write $e_K=v_K(p)$. A statement of existence at absolute ramification index $e$ allows the field $K$ to vary with $e$. The constructions use residue field $\overline{\F}_p$; they make no assertion for every field of that index.

For an elliptic curve, \emph{full rational $p^n$-torsion} means that all of $E[p^n](\overline K)$ lies in $E(K)$. If $E$ is potentially good, its \emph{semistability defect} is the degree of its minimal good-reduction extension after passage to the completed maximal unramified base. The character-recovery problem concerns a Barsotti--Tate group $M/K$ admitting a finite filtration whose graded pieces acquire good reduction. Over a common finite Galois extension where those pieces have good models, inertia acts on their rational covariant Dieudonn\'e modules. The sum of these finite characters, evaluated on elements of order prime to $p$ in the finite inertia quotient, will be denoted $\theta_M$. This is the character defined in Grothendieck's \emph{Cristaux} manuscripts, cote~9, sheet~11.

On sheet~12 Grothendieck asks, for elliptic curves with isomorphic $p$-torsion,
\begin{quote}
``est-il vrai que les fonctions traces $\phi_{E_1}$ et $\phi_{E_2}$ sur le groupe d'inertie [\ldots] sont \'egales ?''
\end{quote}
The manuscript restricts the comparison to elements of order prime to $p$ and explicitly assumes that $K$ contains the $p$th roots of unity. The quotation is from archive sheet~12, which is PDF page~13; the character definition occupies archive sheet~11, PDF page~12 \citep[pp.~11--13]{Grothendieck1969}. We study the refined question of the least higher truncation sufficient for that character, within the stated filtration class.

We use actual finite torsion throughout: an isomorphism $M[p^n]\simeq M'[p^n]$ is an isomorphism of finite group schemes over $K$, equivalently of their finite Galois modules. It is not an isomorphism of rational Tate modules. In every character-index assertion we restrict the base class by $\mu_p\subset K$. Let $\nuMin(p,e,h)$ be the least nonnegative integer $\nu$ for which this character is determined uniformly, among allowed height-$h$ groups over bases of index $e$, by $M[p^{\nu+1}]$. Thus the index $\nu=2$ refers to $p^3$-torsion.

The allowed filtration condition is the monodromy condition defined on archive sheet~10, PDF page~11: a finite composition series of Barsotti--Tate groups whose graded pieces acquire good reduction after finite separable extension. Sheet~11, PDF page~12, defines its sum character \citep[sheets~10--11]{Grothendieck1969}. Proposition~\ref{prop:filtered-non-ht} shows that this condition can hold when the whole Tate representation is not potentially Hodge--Tate. The results below comparing an elliptic curve to its full rational torsion do not by themselves determine $\nuMin$.

\subsection{The complete prime-to-$p$ index}
The first theorem determines the character index at every absolute ramification index prime to $p$, in the full height-two filtration class defined above.

\begin{theorem}[The complete prime-to-$p$ character index]\label{thm:complete-prime-index}
Let $p\ge5$ and let $e$ be a positive multiple of $p-1$ with $p\nmid e$. Consider complete discretely valued characteristic-zero fields with perfect residue field of characteristic $p$, absolute index $e$, and $\mu_p\subset K$. In the height-two Barsotti--Tate filtration class,
\begin{equation}\label{eq:complete-prime-index}
 \nuMin(p,e,2)=\left\lceil\frac{1+\log_p e}{2}\right\rceil.
\end{equation}
Equivalently, $\nuMin(p,p-1,2)=1$, and for every integer $r\ge2$,
\[
 \nuMin(p,e,2)=r\qquad\text{if}\qquad p^{2r-3}<e<p^{2r-1}.
\]
Put $r=\nuMin(p,e,2)$. The upper statement is uniform over every field in the specified base class: $p^{r+1}$-torsion determines the special prime-to-$p$ inertia character. At every such $e$ there are actual allowed groups over some field with residue $\overline{\F}_p$, isomorphic $p^r$-torsion and different special characters. These lower examples are Barsotti--Tate groups; no assertion that every pair comes from elliptic curves is made.
\end{theorem}

The strict endpoints cause no ambiguity because $p\nmid e$. In particular, the complete intervals for indices two, three and four are
\[
 p<e<p^3,\qquad p^3<e<p^5,\qquad p^5<e<p^7,
\]
respectively, with the same divisibility and base-field restrictions. The proof combines an upper comparison theorem with actual lower constructions at every index.

\begin{theorem}[Iterated supersingular recovery]\label{thm:ss-iterated-range}
Let $p\ge5$, $n\ge3$, and let $K$ be a complete discretely valued field of characteristic zero with perfect residue field of characteristic $p$, containing $\mu_p$. Put $e=e_K$ and assume $p\nmid e$. If
\begin{equation}\label{eq:ss-iterated-range}
 e<\frac{(p-1)p^{3n-3}}{p^n-1},
\end{equation}
then any isomorphism $G[p^n]\simeq H[p^n]$ between height-two Barsotti--Tate groups with potentially supersingular reduction implies $\theta_G=\theta_H$. Since $(p-1)\mid e$, the numerical condition is equivalently
\[
 e\le(p-1)\bigl(p^{2n-3}+p^{n-3}\bigr).
\]
\end{theorem}

\begin{corollary}[The full filtration range]\label{cor:complete-prime-recovery}
Under the same base hypotheses, for every $n\ge3$ one has
\[
 e<p^{2n-3}\quad\Longrightarrow\quad\nuMin(p,e,2)\le n-1.
\]
\end{corollary}

Section~\ref{sec:canonical-iteration} proves these statements by repeatedly taking a canonical quotient of one model and the quotient by its transported kernel in the other. Each step loses exactly one torsion level. The graph valuation bound ensures that the first model enters the central range while at least $p^2$-torsion remains, so Theorem~\ref{thm:central-p2-rigidity} applies. The transported kernel need not be canonical in the partner. The sharp mixed obstruction then supplies both the full-class upper threshold and the matching lower examples.

\subsection{Uniform recovery and exact wild bands}
Primitive coordinates give a uniform comparison without a restriction on the $p$-part of the absolute ramification index; a compatible cyclic-point construction gives exact mixed thresholds and infinitely many sharp full-class values.

\begin{theorem}[A uniform range at every index]\label{thm:all-index-upper}
Let $p\ge5$ and $n\ge2$. In the standing height-two filtration class, if
\[
                           e_K\le p^{n-1},
\]
then isomorphic actual $p^n$-torsion forces equality of the special inertia characters. This holds over every standing complete field of that index, with no restriction on $v_p(e_K)$. Equivalently,
\[
 \nuMin(p,e,2)\le\left\lceil\log_p e\right\rceil
\]
for every positive multiple $e$ of $p-1$.
\end{theorem}

\begin{theorem}[The mixed threshold at every index]\label{thm:mixed-all-indices}
Let $p\ge5$, $n\ge2$ and let $e$ be a positive multiple of $p-1$. Put $v=v_p(e)$. Over some standing field of exact index $e$, a whole potentially supersingular height-two group and a group filtered by two potentially good height-one groups of total dimension one have isomorphic actual $p^n$-torsion and different special characters if and only if
\[
                              e>p^{2n-3-v}.
\]
The existence field may always have residue field $\overline{\F}_p$. A negative exponent is interpreted in $\Q_{>0}$.
\end{theorem}

\begin{theorem}[Infinite exact wild bands]\label{thm:wild-bands-exact}
For every $p\ge5$, integer $r\ge0$ and integer $2\le k\le(p^2-1)/2$,
\[
 \nuMin\bigl(p,kp^r(p-1),2\bigr)=
 \begin{cases}
 r+2,&p\nmid k\ \text{or }k=p,\\
 r+3,&p\mid k,\ k\ge2p.
 \end{cases}
\]
In the first case, $p^{r+3}$-torsion recovers the special character over every standing field of that exact index, while some such field supports two actual permitted groups with isomorphic $p^{r+2}$-torsion and different special characters. In the second case, these two levels are $p^{r+4}$ and $p^{r+3}$, respectively. In particular,
\[
 \begin{gathered}
 \nuMin\bigl(p,kp(p-1),2\bigr)=3\quad(2\le k\le2p-1),\\
 \nuMin\bigl(p,p^s(p-1),2\bigr)=s+1\quad(s\ge1).
 \end{gathered}
\]
At $k=2p$, this gives
\[
 \nuMin\bigl(p,2p^{r+1}(p-1),2\bigr)=r+3.
\]
\end{theorem}

Section~\ref{sec:one-slope-all-indices} proves the uniform theorem by comparing two Frobenius chains and the larger exact band by matching residual characters and transporting a canonical quotient. Section~\ref{sec:all-wild-indices} gives a separate graph proof of the sufficient range $e\le(p-1)p^{n-2}$, proves the exact mixed criterion and constructs the lower examples for the bands. The mixed criterion determines exactly when an actual group of one reduction type can have the same truncation as a group of the other type. Comparisons between two supersingular groups outside the proved ranges remain a separate question, so the infinite exact bands do not constitute a formula for every later wild index.


\subsection{Primitive comparisons at arbitrary indices}
The next results distinguish Newton polygons without assuming that the absolute index is prime to $p$. Let $G,H$ be whole potentially supersingular height-two groups over a standing base $K$, with $p\ge5$, $\mu_p\subset K$ and $v=v_p(e_K)$. After completed unramified base change, use the common tame good field $L/K$ of degree $p^2-1$ and normalize $v_L(L^\times)=\Z$. A one-slope model has all nonzero $p$-point values equal to $e_K$. A two-slope model has an integer parameter $z$ with
\[
 0<z<pe_K,\qquad
 (p+1)e_K-z\quad(p-1\text{ points}),\qquad
 z/p\quad(p^2-p\text{ points}).
\]

\begin{theorem}[One-slope recovery]\label{thm:one-slope-all-indices}
If both good models have one first-level Newton slope and
\[
 n\ge v_p(e_K)+2,\qquad G[p^n]\simeq H[p^n],
\]
then $\theta_G=\theta_H$. There is no upper bound on $e_K$.
\end{theorem}

\begin{theorem}[An obstruction to different polygons]\label{thm:mixed-polygon-obstruction}
Let $n\ge1$. If $G[p^n]\simeq H[p^n]$, one model has one slope and the other has two slopes with parameter $z$, then
\begin{equation}\label{eq:mixed-polygon-obstruction}
 v_p(z)\ge\max\{1,2n-1-v_p(e_K)\},\qquad
 e_K>p^{\,2n-2-v_p(e_K)}.
\end{equation}
This assertion does not require unequal special characters.
\end{theorem}

\begin{corollary}[The remaining two-slope comparison]\label{cor:all-index-two-slope-reduction}
Let $p\ge5$, $n\ge2$, and let $M,M'$ belong to the full standing height-two filtration class over $K$. Suppose
\[
 e_K\le p^{\,2n-3-v_p(e_K)},\qquad M[p^n]\simeq M'[p^n].
\]
Their total dimensions and finite determinant factors agree. If their special characters differ, both groups are whole potentially supersingular groups and both good models have two first-level Newton slopes. In every other structural or polygon case the characters agree.
\end{corollary}

Define $\nu_{\mathrm{SS}}(p,e)$ to be the least nonnegative integer $r$ such that actual $p^{r+1}$-torsion recovers the special character uniformly among pairs of whole potentially supersingular height-two groups over the standing fields of index $e$.

\begin{theorem}[The full index and the supersingular index]\label{thm:full-supersingular-maximum}
For every $p\ge5$ and positive multiple $e$ of $p-1$, put $v=v_p(e)$. Then
\[
 \nuMin(p,e,2)=\max\left\{\nu_{\mathrm{SS}}(p,e),
       \left\lceil\frac{1+\log_p e+v}{2}\right\rceil\right\}.
\]
The index on the left concerns the full filtration class. The identity does not evaluate $\nu_{\mathrm{SS}}$ at every wild index.
\end{theorem}

\begin{theorem}[Reduction to a central comparison]\label{thm:all-index-central-reduction}
Under the hypotheses of Corollary~\ref{cor:all-index-two-slope-reduction}, suppose the special characters differ. For the two two-slope models one has
\[
 \max\{v_p(z_1),v_p(z_2)\}\ge v+2,\qquad v=v_p(e_K).
\]
Choose a model $G$ attaining this bound. There are $c$ successive quotients of $G$ by its canonical $p$-lines, paired with quotients of the other model by their transported lines, with
\[
 0\le c\le n-v-2,\qquad r=n-c\ge v+2,
\]
which give an unequal-character actual $p^r$-torsion comparison. The chosen terminal model has parameter $z_c=p^c z_G$ satisfying
\[
 e_K<z_c<pe_K,\qquad
 v_p\bigl((p+1)e_K-z_c\bigr)=v_p(e_K).
\]
The partner is whole potentially supersingular and may have either Newton polygon. The original numerical bound at level $n$ is not asserted at the smaller level $r$.
\end{theorem}

Section~\ref{sec:one-slope-all-indices} proves these assertions. The scaled torsion graph has $2n$ primitive Frobenius dimensions. A one-slope coordinate has a two-dimensional stable Frobenius span, while a two-slope coordinate has a one-dimensional span; this distinction controls overlaps of their nonzero weights. The primitive correspondence is classical \citep[\S\S2.1--2.2, pp.~3--4]{Koch2015}. The two-chain divisibility argument also proves Theorem~\ref{thm:all-index-upper}. At a failed canonical-line matching, Lemma~\ref{os-lem:canonical-matching} supplies an additional truncated valuation-sum inequality. For $v=0$, central $p^2$ rigidity excludes the terminal comparison; for $v>0$, the corresponding assertion with a $p$-divisible high numerator remains open.

\subsection{The first wild character index}
The graph valuation bound also gives an exact full-class answer at the first absolute ramification index divisible by $p$ under the cyclotomic hypothesis.

\begin{theorem}[The general first wild index]\label{thm:first-wild-general}
Let $p\ge5$. For the height-two Barsotti--Tate filtration class over complete discretely valued characteristic-zero fields with perfect residue field, $\mu_p\subset K$ and $e_K=p(p-1)$, one has
\[
                         \nuMin(p,p(p-1),2)=2.
\]
Thus $p^3$-torsion determines the special prime-to-$p$ inertia character over every such field. There are actual unequal-character comparisons at level $p^2$ over a field of that exact index with residue field $\overline{\F}_p$. The lower examples consist of a whole potentially supersingular group and a group filtered by two potentially good height-one groups.
\end{theorem}

Section~\ref{sec:first-wild-general} gives the full proof. For two whole supersingular groups the graph estimate would require
\[
p(p-1)(p^3-1)\ge(p-1)p^4,
\]
whose left side is smaller by $p(p-1)$. The mixed obstruction and determinant comparison settle the other filtration types. An invariant Eisenstein specialization of the cyclic-point construction attains the lower level.


\subsection{The smallest character index}
The next theorem determines the least torsion level at the smallest possible absolute index under the cyclotomic hypothesis.

\begin{theorem}[Minimal cyclotomic ramification]\label{thm:minimal-index}
Let $p\ge5$. For the height-two filtration class over complete discretely valued characteristic-zero fields with perfect residue field, $\mu_p\subset K$ and $e_K=p-1$,
\[
                         \nuMin(p,p-1,2)=1.
\]
Thus $p^2$-torsion determines the special prime-to-$p$ inertia character over every such field. Over $K=\operatorname{Frac}W(\overline{\F}_p)(\zeta_p)$ there are two whole potentially supersingular groups with isomorphic $p$-torsion and different special characters.
\end{theorem}

\begin{corollary}[The exact index-two range]\label{cor:exact-index-two}
Let $p\ge5$ and let $e$ be a positive multiple of $p-1$ with $p\nmid e$. In the same height-two filtration class, with $\mu_p$ in every base field,
\[
              \nuMin(p,e,2)=2\quad\Longleftrightarrow\quad p<e<p^3.
\]
Equivalently $e=t(p-1)$ with $2\le t\le p^2+p+1$ and $p\nmid t$. In particular, at $p=5$ the index is one at $e=4$, and two at every multiple of four from $8$ through $124$ that is prime to five.
\end{corollary}

Section~\ref{sec:minimal-index} proves the theorem by combining the graph valuation bound with an explicit Lubin--Tate twist. The twist changes the characteristic-zero descent trace while Frobenius on $\F_{p^2}$ identifies its first-level representation. The sharp mixed examples then identify both ends of the index-two range. These statements give an exact answer in a Barsotti--Tate class; they do not assert that every lower example is the torsion of an elliptic curve.

\subsection{Central rigidity and prime-to-$p$ recovery}
For the next two theorems, $G$ and $H$ are height-two Barsotti--Tate groups with potentially supersingular reduction, and $p\ge5$. Pass to the completed maximal unramified base, without changing $e=e_K$, and take the common tame good-reduction extension $L/K$ of degree $m=p^2-1$. Normalize $v_L(L^\times)=\Z$. For a two-slope good model, its normalized Hasse value $z$ is the positive integer specified by the nonzero $p$-torsion valuations
\[
 A=(p+1)e-z\quad(p-1\text{ points}),\qquad
 B=z/p\quad(p^2-p\text{ points}),\qquad 0<z<pe.
\]
Equivalently, its deformation Hasse valuation is $h=(p-1)z$. The points of larger valuation form the canonical subgroup $C$. We call the range $e<z<pe$ \emph{central}; it is exactly the range $A<p^2B$ in which every $p$-division point of a nonzero point of $C$ has valuation $A/p^2$. These conventions and their independence of integral parameter changes are justified in Lemma~\ref{rc-lem:newton}.

\begin{theorem}[Central $p^2$ rigidity]\label{thm:central-p2-rigidity}
Let $K$ be a complete discretely valued field of characteristic zero with perfect residue field of characteristic $p\ge5$, and suppose $\mu_p\subset K$. Let $G,H/K$ be height-two Barsotti--Tate groups with potentially supersingular reduction. Suppose the good model of $G$ over the common tame field $L$ has normalized Hasse value $z_G$ satisfying
\[
 e_K<z_G<pe_K,\qquad p\nmid A_G=(p+1)e_K-z_G.
\]
Then every isomorphism $G[p^2]\simeq H[p^2]$ over $K$ implies
$\theta_G=\theta_H$. There is no restriction on $v_p(e_K)$, and no Newton-slope or Hasse-divisibility assumption on $H$.
\end{theorem}

The proof retains two fields attached to a point $Q$ with $pQ\in C\setminus\{0\}$:
\[
 F=L(Q),\qquad F_0=L(q(Q)),\qquad q:G\longrightarrow G/C.
\]
Their degrees over $L$ are $p^2$ and $p$, respectively, and $F/F_0$ is cyclic with exact lower break $(p^2-1)A_G$. A finite torsion isomorphism identifies both fields through their point stabilizers. The lower-break inequality for the transported coordinate then opposes the strict residual-break inequality of a quotient with a $p$-divisible Hasse value. Section~\ref{sec:central-ramification} proves every field-degree, normality and transport assertion used here.

\begin{theorem}[Recovery with one canonical quotient]\label{thm:ss-prime-range}
Let $p\ge5$ and $n\ge3$. Let $K$ be a complete discretely valued field of characteristic zero with perfect residue field of characteristic $p$, containing $\mu_p$, and put $e=e_K$. Suppose $p\nmid e$ and
\begin{equation}\label{eq:ss-prime-range}
 e<\frac{(p-1)p^{2n}}{p^n-1}.
\end{equation}
For any two height-two Barsotti--Tate groups $G,H/K$ with potentially supersingular reduction, an isomorphism $G[p^n]\simeq H[p^n]$ implies $\theta_G=\theta_H$.
\end{theorem}

An invariant augmentation-ideal estimate gives, in a hypothetical unequal-character comparison, the pointwise lower bound $(p-1)/(p^n-1)$ on every nonzero torsion coordinate. For a two-slope model this forces $z/p^{2n-1}\ge(p-1)/(p^n-1)$. Under \eqref{eq:ss-prime-range}, a $p$-divisible Hasse value is either already central or becomes central after one canonical quotient. Theorem~\ref{thm:central-p2-rigidity} then applies. The one-slope case is treated without an upper bound on $e$. This proves Theorem~\ref{thm:ss-prime-range} in Section~\ref{sec:central-ramification}.

\begin{corollary}[Consequences of one canonical quotient]\label{cor:prime-recovery-intro}
For $p\ge5$, a positive multiple $e$ of $p-1$ prime to $p$, and the base class $\mu_p\subset K$, one has
\[
 \nuMin(p,e,2)\le n-1
 \quad\begin{cases}
 n=3,&e<p^3,\\
 n\ge4,&e\le(p-1)(p^n+1).
 \end{cases}
\]
In particular, the following exact values hold:
\begin{align*}
 \nuMin(p,e,2)&=2&&\text{if }p<e<p^3,\\
 \nuMin(p,e,2)&=3&&\text{if }p^3<e\le(p-1)(p^4+1),\\
 \nuMin(p,e,2)&=4&&\text{if }p^5<e\le(p-1)(p^5+1).
\end{align*}
The filtration class and index convention are those fixed above.
\end{corollary}

Theorem~\ref{thm:prime-index-upper} and Corollary~\ref{cor:prime-index-intervals} prove the full-class range $e<p^n$ by recovery of canonical lines. Section~\ref{sec:central-ramification} gives the displayed one-quotient bounds, while Section~\ref{sec:canonical-iteration} proves the complete ranges of Theorem~\ref{thm:complete-prime-index} by iteration. The two comparisons use different ranges of canonical quotient geometry. At $n=3$, the exact mixed obstruction requires the full-class restriction $e<p^3$, even though the supersingular-only range is larger. Section~\ref{sec:primitive-grading} gives a further primitive-function reduction and excludes all three of its numerical configurations by Theorem~\ref{thm:central-p2-rigidity}.

\subsection{Full rational torsion with tame elliptic descent}
The next three results determine exactly which absolute ramification indices admit an additive potentially good elliptic curve with full rational torsion and a prescribed tame defect. The field is allowed to vary; the existence assertions hold with algebraically closed residue field.

\begin{theorem}[Tame full torsion at level one]\label{thm:tame-level-one}
Let $p\ge5$. There exist a complete discretely valued characteristic-zero field $K$ with perfect residue field of characteristic $p$ and $e_K=e$, and an elliptic curve $E/K$ with full rational $p$-torsion, additive potentially good reduction and semistability defect $m$, if and only if
\begin{equation}\label{eq:tame-first-criterion}
 m\in\{2,3,4,6\},\qquad m\mid p+1,\qquad p-1\mid e,
 \qquad e\ge\frac{p^2-1}{m}.
\end{equation}
The residue field may be $\overline{\F}_p$. Consequently the set of possible indices is
\[
 \left\{e\in(p-1)\Z_{>0}:e\ge\frac{p^2-1}{m_{\max}(p)}\right\},
 \qquad
 m_{\max}(p)=
 \begin{cases}
 2,&p\equiv1\pmod{12},\\
 6,&p\equiv5\pmod{12},\\
 4,&p\equiv7\pmod{12},\\
 6,&p\equiv11\pmod{12}.
 \end{cases}
\]
For $p=3$ and a specified tame defect, $m=2$ occurs exactly at even indices $e\ge4$, and $m=4$ occurs at every positive even index. This last assertion concerns tame inertia only.
\end{theorem}

\begin{theorem}[Every tame defect and every level]\label{thm:tame-full-level}
Fix $p\ge5$, $n\ge1$, and $m\in\{2,3,4,6\}$ with $m\mid p+1$. Let $e$ be a positive multiple of $p-1$, and put $s=v_p(e/(p-1))$. There exist a complete discretely valued field $K$ with $e_K=e$ and perfect residue field of characteristic $p$, and an elliptic curve over $K$ with full rational $p^n$-torsion and exact semistability defect $m$, if and only if
\begin{equation}\label{eq:tame-full-criterion}
 s\ge n-1,\qquad
 e\ge\frac{p^2-1}{m}\,
       p^{\max\{2n-2,\,4n-4-s\}}.
\end{equation}
Every permitted index is attained with residue field $\overline{\F}_p$.
The same assertion holds for $p=3$ and specified tame defect $m=2$ or $4$. The wild characteristic-three defects are treated separately in Theorem~\ref{thm:three-range}.
\end{theorem}

\begin{corollary}[The least level forcing good reduction]\label{cor:tame-forcing-level}
Let $p\ge5$ and $e\in(p-1)\Z_{>0}$. If $m_{\max}(p)e<p^2-1$, put $N=0$. Otherwise set
\[
 s=v_p(e/(p-1)),\qquad
 L=\left\lfloor\log_p\frac{m_{\max}(p)e}{p^2-1}\right\rfloor,
 \qquad
 N=1+\min\left\{s,\left\lfloor\frac L2\right\rfloor,
                         \left\lfloor\frac{L+s}{4}\right\rfloor\right\}.
\]
Among potentially good elliptic curves over all complete fields of absolute index $e$ with perfect residue field of characteristic $p$, the least positive level $n$ at which full rational $p^n$-torsion forces good reduction is $N+1$. For $p=3$, the same formula with $m_{\max}=4$ gives the least level within the class with tame inertia image.
\end{corollary}

Section~\ref{sec:tame-full-torsion} proves necessity from the smallest and largest exact-order formal torsion valuations. For sufficiency it uses the regular full-level deformation surface: a tame action scales both basis-point parameters, a divisorial valuation prescribes their extrema, and an invariant transverse section gives an actual elliptic curve over a complete field of exactly the required index. This last step restores an algebraically closed residue field after an intermediate valuation with residue field $\overline{\F}_p(z)$.

The basic full-division-field bound
\[
 D=(p^2-1)p^{2n-2}
\]
is already the lower bound in \citet[Theorem~1.1, p.~2 of the author version]{Smith2023}, in the setting of a good supersingular elliptic curve at a prime of a number field. Our local proof includes arbitrary perfect residue fields. The assertions above further specify the exact tame defect, the dependence on $s=v_p(e/(p-1))$, and realization at every permissible index; no priority claim is made for the basic bound or the torsion-valuation classification.


\subsection{Wild full torsion in residue characteristic three}
The following theorem treats the wild defects at $p=3$. Together with Theorem~\ref{thm:tame-full-level}, it gives the complete existential answer for additive elliptic curves in this characteristic.

\begin{theorem}\label{thm:three-range}
Let $e=2u$ be a positive even integer, put $s=v_3(u)$, and let $n\ge1$. For each $m\in\{3,6,12\}$, an elliptic curve with additive reduction, exact semistability defect $m$, and full rational $3^n$-torsion exists over some complete discretely valued field of absolute ramification $e$ and perfect residue field if and only if
\begin{equation}\label{eq:wild-criterion}
 s\ge n-1,\qquad
 me\ge 8\,3^{\max\{2n-2,\,4n-4-s\}}.
\end{equation}
The field can be chosen with residue field $\overline{\F}_3$.

Allowing every nontrivial defect, the largest possible level is
\begin{equation}\label{eq:maximum-three}
 N(e)=1+\min\left\{s,\left\lfloor\frac{L+1}{2}\right\rfloor,
                         \left\lfloor\frac{L+s+1}{4}\right\rfloor\right\},
 \qquad L=\lfloor\log_3u\rfloor.
\end{equation}
Among potentially good elliptic curves, full rational $3^{N(e)+1}$-torsion forces good reduction, and that level is least.
\end{theorem}

The theorem describes how much ramification is needed to hide additive reduction from full rational torsion. The smallest absolute index admitting full rational $27$-torsion with additive reduction is $162$. The index $216$ also admits such a curve, of defect twelve; its significance is that the tame construction does not attain that level there. We give the parameters of both examples in Section~\ref{sec:three}.

The proof combines two ingredients. First, rational full level structure gives an exact sum of the valuations of primitive torsion coordinates. A wild automorphism forces the largest such valuation to be divisible by three. Second, one must realize the numerical possibility by a deformation carrying the complete descent action. Regularity of a level ring alone does not provide this realization. We construct an invariant equation on a weighted chart and prove that it has a unique simple branch; that branch is the required discrete valuation ring.

\subsection{Ordinary reduction and boundary examples}
The following results distinguish unramified torsion of an ordinary abelian variety from the broader filtration condition used in the character problem.

\begin{theorem}[Ordinary reduction from unramified $p$-torsion]\label{thm:ordinary-unramified-p}
Let $K$ be complete discretely valued of characteristic zero with perfect residue field of odd characteristic $p$. Let $A/K$ be an abelian variety of any dimension with potentially good ordinary reduction. If $A[p]$ is unramified, then $A$ has good reduction. Neither $\mu_p\subset K$ nor a condition on the degree of a polarization is required.
\end{theorem}

\begin{corollary}[Defect two with full rational $p$-torsion]\label{cor:defect-two-full-p}
For every odd prime $p$, there are a complete field $K$ with residue field $\overline{\F}_p$, absolute ramification index $(p^2-1)/2$ and $\mu_p\subset K$, and an elliptic curve $E/K$ with full rational $p$-torsion and additive potentially good supersingular reduction of exact semistability defect two. Its special inertia involution has trace $-2$.
\end{corollary}

\begin{proposition}[An explicit wild cubic]\label{prop:explicit-cubic-defect-six}
Over $K=\Q_3(\zeta_3)$, the elliptic curve
\[
                         E:\ y^2+y=x^3
\]
has full rational $3$-torsion, additive potentially supersingular reduction and exact semistability defect six, also after unramified base change. Its finite special inertia image is cyclic of order six, and its involution has trace $-2$.
\end{proposition}

\begin{proposition}[Good graded pieces without potential Hodge--Tateness]\label{prop:filtered-non-ht}
For every odd $p$ and $K=\Q_p(\zeta_p)$, there is a height-two Barsotti--Tate group $G/K$ in an exact sequence
\[
0\longrightarrow(\Q_p/\Z_p)_K\longrightarrow G
 \longrightarrow(\Q_p/\Z_p)_K\longrightarrow0
\]
whose rational Tate representation is not potentially Hodge--Tate and hence not potentially semistable. The special character defined by its two good graded pieces is the constant character $2$.
\end{proposition}

Section~\ref{sec:early-reduction} proves these assertions. Corollary~\ref{cor:defect-two-full-p} is the defect-two specialization of Theorem~\ref{thm:tame-level-one}; Proposition~\ref{prop:explicit-cubic-defect-six} is proved by an explicit good model and its faithful descent action. The representation in Proposition~\ref{prop:filtered-non-ht} is the classical logarithmic cyclotomic example displayed by \citet[Section~II.1.2, p.~9 of the 2002 author version]{Berger2004}. The proposition constructs its actual Barsotti--Tate group and identifies what it says about the filtration class; it makes no new claim in Sen theory.

\subsection{Height one and full first-level elliptic torsion}
The first layer already recovers the height-one character, while elliptic full torsion has additional constraints that depend on the particular base field.

\begin{proposition}[The height-one index]\label{prop:height-one-index}
Let $p$ be any prime and let $K$ satisfy the standing complete-field hypotheses, with $\mu_p\subset K$. For every height-one group $M$ in the allowed Barsotti--Tate filtration class, its special character on prime-to-$p$ inertia is determined by $M[p]$. Thus
\[
                         \nuMin(p,e_K,1)=0.
\]
This includes $p=2$.
\end{proposition}

\begin{theorem}[A full-three-torsion curve of defect twelve]\label{thm:explicit-hesse-twelve}
Put
\[
 K=\Q_3(i,\zeta),\qquad i^2=-1,\quad \zeta^2+\zeta+1=0,
 \qquad \pi=\zeta-1,\qquad t=1+i\pi.
\]
Then $e_K=2$ and the residue field is $\F_9$. The elliptic curve
\[
 E:\ y^2=x^3-27t(t^3+8)x+54(t^6-20t^3-8)
\]
has full rational $3$-torsion and additive potentially supersingular reduction of exact semistability defect twelve. The same exact defect holds after unramified base change.
\end{theorem}

\begin{theorem}[A variable Hesse defect]\label{thm:hesse-variable-defect}
Let $K$ be a complete discretely valued field of characteristic zero with algebraically closed residue field of characteristic three, containing $\mu_3$ and of absolute ramification index $e$. Choose a uniformizer $\varpi$ and an integer $r$ such that
\[
                         0<2r<e,\qquad 3\nmid e+r.
\]
Put $t=1+\varpi^r$. The elliptic curve
\[
 E_t:\quad y^2=x^3-27t(t^3+8)x+54(t^6-20t^3-8)
\]
has full rational $3$-torsion and potentially supersingular reduction of exact semistability defect
\[
                             m=\frac{12}{\gcd(4,e+r)}.
\]
There is an explicit good-reduction extension of this degree, and the ramification degree of every good-reduction extension is divisible by $m$.
\end{theorem}

\begin{corollary}[Every first-level defect in characteristic three]\label{cor:three-first-level-table}
Allow the complete field to vary, with arbitrary perfect residue field of characteristic three. An additive potentially good elliptic curve with full rational $3$-torsion and specified defect $m$ exists at exactly the indices in the following table:
\[
\begin{array}{c|c}
 m&\text{absolute ramification indices}\\\hline
 2&e\in2\Z,\ e\ge4\\
 3&e\in2\Z,\ e\ge4\\
 4&e\in2\Z,\ e\ge2\\
 6&e\in2\Z,\ e\ge2\\
 12&e\in2\Z,\ e\ge2.
\end{array}
\]
There are no other nontrivial defects. In particular every positive even index is attained. Over the particular field $\Q_3(\zeta_3)$ there is an example of defect six.
\end{corollary}

\begin{corollary}[The cyclotomic base field]\label{cor:cyclotomic-full-p}
For $p\ge5$, no elliptic curve defined over $\Q_p(\zeta_p)$ has both full rational $p$-torsion and additive potentially good reduction. For $p=3$ such a curve exists.
\end{corollary}

Section~\ref{sec:first-level-refinements} proves these statements. The Hesse family and its preservation of symplectic $3$-torsion structures are classical; see \citet[\S1, Theorem~1.1, p.~2]{AnemaTopTuijp2018}, over perfect fields of characteristic different from three. Here the generic field has characteristic zero. The proofs supply the points and the exact local good-reduction fields for the displayed curves. The first-level table is a direct specialization of the complete tame and wild existence criteria. The fixed cyclotomic field requires an additional residue-field argument at $p=5$: an existential example at the same absolute index over a different field would not settle that corollary. The classifications of \citet{OzekiYoshida2026} concern curves defined over $\Q_p$; they are neighboring context and are not used as a proof of this assertion for curves defined over $\Q_p(\zeta_p)$.


\subsection{A cyclic point and a cyclotomic obstruction}
The next results concern inputs to the comparison problem. They are stated here in their independent field-theoretic and formal-group forms.

\begin{theorem}\label{thm:cyclic-bound}
Let $p\ge5$, let $R$ be a complete discrete valuation ring with algebraically closed residue field of characteristic $p$, and put $a=v_R(p)$ and $v=v_p(a)$. If a one-dimensional formal group of special height two has an $R$-rational point of exact order $p^n$, with $n\ge2$, then
\begin{equation}\label{eq:cyclic-bound}
 a>(p^2-1)p^{2n-3-v}.
\end{equation}
A negative exponent on the right is interpreted in $\Q_{>0}$.
\end{theorem}

This inequality records the number of times a rational inverse chain can follow the larger Newton slope. Each such step consumes one factor of $p$ in $a$. The valuations themselves are the canonical-subgroup valuations described by \citet[Theorem~4.4 and Corollaries~4.5--4.6, pp.~10--11 of the author version]{Smith2023}. Section~\ref{sec:cyclic} gives the argument for arbitrary height-two formal groups and derives the strict inequality, including its equality boundary.

\begin{theorem}\label{thm:cyclotomic-exact}
Let $p\ge3$, and suppose that $K$ has algebraically closed residue field, $e_K=p(p-1)$, $\mu_p\subset K$, and $\mu_{p^2}\not\subset K$. Choose a Cohen coefficient field $K_0\subset K$, put $K_1=K_0(\mu_p)$, and write $d=v_K(\mathfrak D_{K/K_1})$. The last upper ramification break of $K(\mu_{p^3})/K$ is exactly
\begin{equation}\label{eq:exact-second-break}
 \delta_2=(p-1)(2p+1)-d\ge p(p-1).
\end{equation}
\end{theorem}

This identity follows by taking the normal closure of the degree-$p$ extension $K/K_1$ and applying Herbrand's theorem. It supplies the determinant obstruction in the formal comparison proofs. A separate argument with differents gives the following bound, without computing that normal closure.

\begin{theorem}\label{thm:cyclotomic-bound}
Let $p\ge3$. Suppose that $K$ has algebraically closed residue field, $e_K=p(p-1)$, $\mu_p\subset K$, and $\mu_{p^2}\not\subset K$. Then $C=K(\mu_{p^3})$ is cyclic of degree $p^2$ over $K$. If $\delta_1<\delta_2$ are its upper ramification breaks, in the classical convention without a conductor shift, then
\begin{equation}\label{eq:cyclotomic-bound}
 \delta_1+p\delta_2\ge p(p^2-p-1).
\end{equation}
Consequently $\delta_1\le p-1$ implies
\begin{equation}\label{eq:second-break}
 \delta_2\ge p^2-p-2+\frac1p.
\end{equation}
\end{theorem}

The inequality follows by comparing differents in two towers over a Cohen field. Both arguments apply to infinite perfect residue fields. A proposed common torsion field has to contain the cyclotomic extension through its determinant; an upper bound on the torsion field's last break can therefore contradict~\eqref{eq:exact-second-break}.

\subsection{Sharp comparison thresholds}
The following special cases illustrate the character thresholds at small ramification indices. Throughout the following statement, the class defining $\nuMin$ includes the potentially good graded filtrations of Section~1.

\begin{theorem}\label{thm:thresholds-intro}
Let $p\ge5$ and assume $\mu_p\subset K$.
\begin{enumerate}
\item For $n\ge2$ and a positive multiple $e$ of $p-1$ prime to $p$, a potentially supersingular height-two group and a group filtered by two potentially good height-one groups of total dimension one can have isomorphic $p^n$-torsion and different special characters over some field of index $e$ if and only if $e>p^{2n-3}$.
\item Among indices $e$ divisible by $p(p-1)$ with $v_p(e)=1$, such a comparison at level $p^3$ exists if and only if $e>p^2$.
\item Over all complete fields with perfect residue field of absolute index $2p(p-1)$, the least uniform character index is
\[
                         \nuMin(p,2p(p-1),2)=3.
\]
\item At the first wild index, $\nuMin(p,p(p-1),2)=2$ for every $p\ge5$.
\item If $p\nmid e$ and $p^3<e<p^4$, then $\nuMin(p,e,2)=3$.
\end{enumerate}
\end{theorem}

Part~(1) is Corollary~\ref{cor:sharp-mixed-prime}, part~(2) is Corollary~\ref{cor:sharp-mixed-wild}, part~(3) is Theorem~\ref{thm:second-wild-exact}, part~(4) is Theorem~\ref{thm:first-wild-general}, and part~(5) is Corollary~\ref{cor:prime-index-intervals}. The existence assertions allow the base to vary; part~(3) is a uniform recovery assertion. Their proofs explain why a cyclic point can conceal a special character: quotienting by the point produces a finite extension of rank-one characters, and Kummer theory lifts exactly that extension to a filtered Barsotti--Tate group. For the upper bound, the common determinant forces the cotangent weights of the graph quotient to agree modulo $p-1$. This strengthens the coefficient estimate in Bondarko's integral kernel argument.

\subsection{The first unresolved wild coefficient}
The next theorem narrows the first coefficient beyond the exact wild bands to one necessary configuration of actual torsion representations.

\begin{theorem}[The first coefficient beyond the exact bands]
\label{thm:first-boundary-normal-form}
Let $p\ge5$ and $r\ge1$, and put
\[
 k_0=\frac{p^2+1}{2},\qquad b_0=\frac{p-1}{2},\qquad
 e=k_0(p-1)p^r,\qquad n=r+3.
\]
Suppose that two groups in the standing height-two filtration class
over a field of absolute ramification index $e$ have isomorphic
actual $p^n$-torsion and different special inertia characters.
Then both are whole potentially supersingular groups with two-slope
good models. In the common tame normalization their Hasse parameters
and special tangent scalars satisfy
\[
 z_1=z_2=b_0p^n,\qquad \frac{\alpha_2}{\alpha_1}=-1.
\]
Both residual $\mathbf F_p$-representations over $K$ are split.
Their canonical and complementary lines have distinct inertia
characters, and the given isomorphism sends each canonical line
to the complementary line of the other representation.
\end{theorem}

This is a necessary form for an assumed unequal-character comparison.
It does not assert that this form occurs or determine the character
index at $e=k_0(p-1)p^r$.

Section~\ref{sec:first-boundary-normal-form} proves the result using the general two-chain congruence and the actual quotient markings.

\begin{corollary}[Point fields at the first boundary]
\label{cor:boundary-point-field-degrees}
Under the assumptions of
Theorem~\ref{thm:first-boundary-normal-form}, pass to the completed
maximal unramified base and the common tame good field $L$.
For every exact $p^j$-point $Q$ of either group whose $p^{j-1}$-multiple is a
nonzero canonical point, write $[L(Q):L]=p^d$. Then
\[
 d\ge\max\{0,2j-1-r\}\quad(2\le j<n),\qquad
 d\ge r+5\quad(j=n).
\]
Each relative step at levels $r+1,r+2,r+3$ has degree $p^2$.
In particular every canonical $p^{n-1}$-point field has degree
at least $p^{r+3}$. If $r=1$ and $n=4$, the degrees at levels
$p^2,p^3,p^4$ are respectively
\[
 p^2,\qquad p^4,\qquad p^6.
\]
\end{corollary}

Proposition~\ref{prop:proper-canonical-point-fields} proves a
more general proper-level bound for an exchanged central
comparison. The point-field restrictions above do not settle
whether the boundary comparison occurs.

\begin{corollary}[Nonsplit Cartan exclusion at the first boundary]
\label{cor:boundary-nonsplit-cartan-image}
Let $p\ge5$, $r\ge1$, $n=r+3$, and
\[
                    e=\frac{p^2+1}{2}(p-1)p^r.
\]
Suppose that two groups in the standing height-two filtration class
over a field of absolute index $e$ have isomorphic actual
$p^n$-torsion and different special inertia characters.
Pass to the completed maximal unramified base and its common
tame good extension $L$ of degree $p^2-1$.
For each group, the finite image of
$\operatorname{Gal}(\overline L/L)$ on its $p^n$-torsion is not
contained in any conjugate of
\[
                    (1+p\mathcal O)/(1+p^n\mathcal O),
\]
where $\mathcal O$ is the unramified quadratic
$\mathbf Z_p$-algebra and $\mathcal O/p^n\mathcal O$ acts on
$(\mathbf Z/p^n\mathbf Z)^2$ through its regular representation.
\end{corollary}

This is a necessary restriction on an assumed actual comparison.
It does not exclude the remaining non-Cartan images or determine
the boundary character index. Section~\ref{sec:boundary-nonsplit-image-proof}
proves the restriction by constructing an actual invariant-ring map
under the additional image-containment assumption.

\begin{corollary}[The irreducible quadratic self-twist case]
\label{cor:boundary-nonsplit-self-twist}
Under the assumptions of
Corollary~\ref{cor:boundary-nonsplit-cartan-image}, the finite
representation of either group admits no tame quadratic self-twist
with an invertible intertwiner whose reduction has irreducible
quadratic characteristic polynomial over $\mathbf F_p$.
\end{corollary}


The proof in Section~\ref{sec:boundary-nonsplit-image-proof} covers
arbitrary perfect residue fields. The split-intertwiner case is not
part of this corollary.


\subsection{Relation to earlier work and scope}
Drinfeld's regularity and level-parameter theorems are the foundation of the full-torsion constructions. For the tame criterion we use the original \citet[Propositions~4.2--4.3, pp.~570--573]{Drinfeld1974}; the primitive factorization used in the wild characteristic-three proof is stated in \citet[Theorem~2.1.2, p.~437; Proposition~2.2.3, pp.~439--441]{CarterStrauch2021}. The equivalence between elliptic and Barsotti--Tate deformations is the Serre--Tate theorem, in the form of \citet[Theorem~1.2.1]{Katz1981}. The divisibility of a discriminant for full rational $3$-torsion follows from the Hesse-pencil theorem of \citet[Theorem~1.1, p.~2; Proposition~5.1, pp.~4--5]{AnemaTopTuijp2018}.

Proposition~\ref{prop:divisorial-full-cartan} computes the full principal nonsplit Cartan quotient of the universal height-two deformation at every level $p^n$, for odd $p$ and $n\ge2$, with monomial weights $(p^2+p+1,p^2)$. Its ramification index and purely inseparable residue degree are both $p^{n-1}$, and its graded fixed field is $k(X^{p^{n-1}},Y^{p^{n-1}})$. Proposition~\ref{prop:divisorial-last-cartan} gives the separate last-layer invariants, while Proposition~\ref{prop:divisorial-scalar-level-four} computes the distinct full scalar subgroup at level four. At $p=5,n=4$, a further tame quotient of ramification index $12$ with unchanged residue field has absolute index exactly $6500$ (Corollary~\ref{cor:divisorial-full-cartan-6500}). This statement concerns the specified divisorial valuation with imperfect residue field; it does not identify its curve specializations with arbitrary boundary comparisons.

The local rings of this Cartan construction have a stronger property. Put $\pi=x^{1-p}y^p$, $t=x^{p^2}/y^{p^2+p+1}$ and $B_c=(R_n[\pi,t,t^{-1}])_{(\pi,t-[c])}$ for $c\in k^*$. Proposition~\ref{prop:cartan-torus-regularity} proves that $B_c^C$ and its reduced vertical divisor are regular, with divisor ring $k[t^{p^{n-1}},t^{-p^{n-1}}]_{(t^{p^{n-1}}-c^{p^{n-1}})}$. Its local equation $\rho$ satisfies $p=u\rho^{(p^4-1)p^{n-1}}$ for a unit $u$. Thus every local mixed-characteristic specialization of these charts has absolute index divisible by $(p^4-1)p^{n-1}$. At $p=5,n=4$, this is $78000$; an additional field extension of ramification index at most $12$ gives base index at least $6500$. The proof uses the principal-augmentation criterion of \citet[Theorem~2 and Proposition~5, pp.~64--69]{KiralyLutkebohmert2013}. Corollary~\ref{cor:boundary-nonsplit-cartan-image} constructs the invariant-ring map under an actual nonsplit-Cartan image hypothesis and contradicts this divisibility for every first-boundary index. Lemma~\ref{lem:scalar-compatible-frame} supplies the coefficient-compatible universal frame over arbitrary algebraically closed residue extensions; its proof derives the required scalar extension degree by degree from \citet[Propositions~4.1--4.3, pp.~570--573]{Drinfeld1974}. If the torus residue is transcendental over $\overline{\mathbf F}_p$, the map from a closed chart need not be local, but the exact unit identity still gives the contradiction. The image hypothesis is not forced for every surviving comparison.

Proposition~\ref{prop:cartan-orbit-data} shows that the maximal primitive-point orbit sizes do not bound the full image by a Cartan: determinant-constrained congruence groups have the same orbits and are noncommutative. Lemma~\ref{lem:quadratic-self-twist-cartan} gives the stronger Cartan containment supplied by an actual quadratic self-twist, with a distinct split case. Corollary~\ref{cor:boundary-nonsplit-self-twist} excludes its irreducible-intertwiner case for an actual first-boundary comparison.

The local divisor class is an additional constraint on specialization. Lemma~\ref{lem:cartier-index-contact} shows that a single vertical multiplicity $N$ and Cartier index $i$ force $N/\gcd(N,e_K)\mid i$ along a local mixed-characteristic curve. Example~\ref{ex:noncartier-quotient-specialization} applies the known cyclic quotient of \citet[Example~5.10]{LiedtkeMatsumoto2018} to show why normality and generic \'etaleness alone do not preserve a divisorial index bound. Proposition~\ref{prop:boundary-specialization-contact} separately forces contact divisible by $5$ for a specified upstairs chart carrying an actual boundary comparison. Example~\ref{ex:fierce-transverse-specialization} shows how a downstairs transverse curve may acquire such contact.

For finite flat group schemes, \citet[Theorems~A--B, p.~661]{Bondarko2006} proves an almost-fullness theorem for the generic-fibre functor and finite torsion criteria for good reduction. These are uniform extension results; they do not on their own give an exact least character index in height two. \citet[Theorem~1.1, p.~1 of the preprint]{Hattori2008} determines tame characters of ramification graded pieces of a finite flat group scheme by its Abbes--Saito filtration. That theorem concerns generic finite representations and their ramification; the character here is the finite descent character of good special graded pieces. Recovering the latter from a given truncation requires an additional argument.

The full-division-field lower bound and higher-torsion valuation classification of \citet[Theorem~1.1, p.~2; Theorem~4.4 and Corollaries~4.5--4.6, pp.~10--11 of the author version]{Smith2023}, and the reduction restrictions for a rational point of order $p^n$ in \citet[Theorems~1.1--1.3, pp.~1--2]{Melistas2022}, address closely related local questions. The numerical lower bound $(p^2-1)p^{2n-2}$ is therefore credited to Smith. The mixed criterion in Theorem~\ref{thm:mixed-all-indices} additionally constructs actual groups of two different reduction types at every admissible absolute index, using a full integral descent action. The exact fixed-defect, every-index assertions~\eqref{eq:tame-full-criterion} and~\eqref{eq:wild-criterion} additionally require an equivariant elliptic realization and the refined valuation of the base index. More recently, \citet{OzekiYoshida2026} determine first-level point fields and ramification breaks for elliptic curves over $\Q_p$; its Section~3.4 also treats bounds for higher torsion. Those results concern fixed elliptic curves over $p$-adic fields, whereas our existence theorem varies the base and permits an algebraically closed residue field.

At the level of entire Galois images, \citet[Theorem~1.3, p.~2]{BisattEtAl2026} classify the $p$-adic images of elliptic curves over $\Q_p$, for $p>7$, whose mod-$p$ image lies in a nonsplit Cartan normalizer; in the higher semistability-defect case their conclusion has an explicit additional bound relating $p$ to the valuation of the $j$-invariant. Their fixed base field and residual-image hypothesis differ from the torsion-character recovery problem over the ramified bases considered here.

The Newton inputs to Theorem~\ref{thm:central-p2-rigidity} agree, for elliptic formal groups, with the canonical-subgroup valuations of \citet[Theorem~4.4 and Corollaries~4.5--4.6]{Smith2023}; Lemma~\ref{rc-lem:newton} derives the required formulas for height-two formal models. The elementary lower-ramification estimate in Lemma~\ref{cr-lem:valuation} belongs to the classical framework of \citet[Chapter~IV, \S1]{Serre1979}. The comparison here retains the cyclic extension $L(Q)/L(q(Q))$, its exact break, and its identification under a specified $p^2$-torsion isomorphism. In contrast, the first-level classifications of \citet[\S\S2.1, 2.3]{OzekiYoshida2026} concern point fields of elliptic curves over $\Q_p$, while their higher-level Proposition~3.25 bounds the last upper break of the full division field. Those stated results do not provide the transported relative-field comparison proved here. Hattori's theorem concerns the generic ramification-graded characters described above. These are comparisons of the precise statements inspected in a bounded primary-literature search, not a claim that the relative-field method or this criterion has no earlier occurrence.

All conversions between upper and lower numbering, and all different calculations, use the classical conventions of \citet[Chapters~III--IV]{Serre1979}. The full-class prime-to-$p$ table is complete. Theorems~\ref{thm:mixed-all-indices} and~\ref{thm:wild-bands-exact} give the complete mixed criterion and infinite exact wild bands. The final section identifies the remaining full-class questions outside those bands and the distinct supersingular-only and fixed-field problems. No claim of priority is made.

Recent work also studies ramification through integral lattices and actual point fields. \citet[Theorem~11.13, p.~43]{DembeleEtAl2026} bound the residual point-field breaks of nonordinary rank-two crystalline representations over ramified quadratic $2$-adic fields. The global theorem of \citet[Theorem~1.1, p.~1]{OpenAI2026} concerns modularity of odd regular two-dimensional $2$-adic representations. Their base fields, torsion levels, and conclusions differ from the odd-prime character thresholds and exact elliptic full-torsion criteria considered here. We use neither result as a proof input.

For the supersingular comparison we also prove the following elementary obstruction in Proposition~\ref{fc-thm:valuation}: if two height-two potentially supersingular groups over \(K\) have isomorphic \(p^n\)-torsion and unequal special inertia traces, then
\[
e_K(p^n-1)\geq(p-1)p^{2n-2}.
\]
No condition on \(v_p(e_K)\) is needed for this inequality. At \(n=3,e_K=p(p-1)\), its left side falls short of its right side by \(p(p-1)\), proving the supersingular upper step of Theorem~\ref{thm:first-wild-general} for every \(p\ge5\).
